% ==================== 二、软件的主要功能 ====================
\section{软件的主要功能}

\subsection{四类角色与各自的使用场景}

系统只有四类角色，且每类角色只在一端活动——这是 REQ-001「最小权限」在界面层面的直接体现：不同角色看到的不是同一套菜单的裁剪版，而是四套各自独立的前端。

\begin{table}[H]
\centering
\caption{角色、使用端与核心场景}
\tabfont
\begin{tabularx}{\textwidth}{@{}L{2.6cm}L{2.8cm}X@{}}
\toprule
角色 & 使用端 & 核心场景 \\
\midrule
顾客 & 手机 H5 & 搜房 $\rightarrow$ 看房型与实时余量 $\rightarrow$ 下单 $\rightarrow$ 15 分钟内支付 $\rightarrow$ 到店入住 $\rightarrow$ 退房后评价；入住日 18:00 前可自助取消 \\
前台 & 桌面工作台 & 看房态看板与今日待办 $\rightarrow$ 排房 $\rightarrow$ 3 步内办完入住 $\rightarrow$ 退房与派打扫 $\rightarrow$ 线下开单 $\rightarrow$ 18:00 后替顾客提取消申请 \\
店长 & 桌面后台 & 按日历定房价与线下保留额度 $\rightarrow$ 审批改价 / 取消 / 超限开单 / 维修派单 $\rightarrow$ 排班 $\rightarrow$ 看评价与绩效；对房态\textbf{只读} \\
平台地区负责人 & 桌面后台 & 向门店下发价格区间 $\rightarrow$ 审核跨店绩效 $\rightarrow$ 选派与调动店长 $\rightarrow$ 审批删除原始数据申请；\textbf{不得}代门店办理入住或改房态 \\
\bottomrule
\end{tabularx}
\end{table}

\subsection{九个业务模块与一个支撑模块}

模块按\textbf{业务能力}切分，不按角色切分：同一个模块的接口会同时被两端调用（例如「库存」既被顾客下单扣减，也被前台线下开单扣减、被店长改上限）。这样切的好处是一个业务概念只有一处实现，代价是模块数量与角色数量不对应，需要靠权限模块在入口处拦截。

\begin{table}[H]
\centering
\caption{模块与主要功能}
\tabfont
\begin{tabularx}{\textwidth}{@{}L{2.9cm}X@{}}
\toprule
模块 & 主要功能 \\
\midrule
M1 账号与权限 & 用户与角色维护；登录认证与会话签发；每个请求校验「角色 + 数据范围」；员工账号由上级分配；店长调动时权限实时回收 \\
M2 门店与房源 & 门店档案与营业状态；房型档案与上下架；房间档案（房号、楼层、房型归属、当前状态、预计修复日期）；对外提供房间总数与维修间数 \\
M3 定价 & 按「房型 $\times$ 日期」维护日价与价格来源；平台价格区间的接收与校验；改价申请与审批；为下单提供带锁取价 \\
M4 库存 & 按「房型 $\times$ 日期」逐晚维护上限、线上已售、线下保留、线下已售；下单条件扣减、取消与超时释放；按公式重算线上可售上限 \\
M5 预订与订单 & 订单状态机；下单（建单 + 扣库存同事务）；改期改房型；取消判定与退款发起；22:00 未到店与次日 12:00 逾期处理；订单查询 \\
M6 支付与退款 & 创建支付单；回调验签、按支付单号判重；原路退款；发布「支付成功」事件；每日对账与补单 \\
M7 房态与工单 & 房间状态机；空闲房间查询与排房占房；办理入住与退房；生成打扫与维修工单；工单执行与关闭 \\
M8 评价与绩效 & 已退房订单的一次性评价（1--5 星 + 文字）；门店与店长绩效汇总；跨店绩效供平台审核；排班与打卡 \\
M9 审计与数据保留 & 操作日志写入与分角色查询；3 年保留策略；半年绩效归档；删除原始数据的审批与执行 \\
M10 通知 & 订阅订单、支付、工单事件，按模板发短信与站内消息；失败重试与降级；失败不触发订单回滚 \\
\bottomrule
\end{tabularx}
\end{table}

其中 M1 与 M9 是\textbf{横切模块}——每一个业务模块都要校验权限、都要写审计日志，因此它们在组件图上不画十根箭头，而是合并成一根「公共支撑」箭头（见 3.3 节）。

\subsection{需求到模块的追溯矩阵}

下表把实验一的 36 条需求逐条映射到模块。功能需求 29 条中 28 条落到具体模块；非功能需求与范围约束横切全系统，在架构设计一节展开；两条待确认需求明确标注为不覆盖。

{\tabfont
\begin{xltabular}{\textwidth}{@{}p{1.55cm}p{4.15cm}p{3.5cm}X@{}}
\caption{REQ—模块追溯矩阵}\label{tab:trace}\\
\toprule
编号 & 需求名称 & 主要模块 & 落点说明 \\
\midrule
\endfirsthead
\toprule
编号 & 需求名称 & 主要模块 & 落点说明 \\
\midrule
\endhead
\midrule
\multicolumn{4}{r}{\zihao{6}续下页}\\
\endfoot
\bottomrule
\endlastfoot
REQ-001 & 角色定义与最小权限 & M1 & 角色与数据范围在会话中签发，网关逐请求校验 \\
REQ-002 & 顾客查询与预订 & M2、M4、M5 & 房源提供可售房型，库存给出余量，订单落单 \\
REQ-003 & 退房与评价提交 & M7、M8 & 退房改订单与房态；评价限已退房且一次 \\
REQ-004 & 前台办理入住 & M7 & 排房 + 入住两动作，同事务改房间与订单 \\
REQ-005 & 房间状态机 & M7 & 六态转换表驱动，非法转换抛异常 \\
REQ-006 & 店长门店管理 & M2、M3、M4、M7 & 日价在 M3，保留与上限在 M4，营业状态在 M2，房态只读在 M7 \\
REQ-007 & 平台地区负责人管理 & M1、M3、M8 & 调动在 M1，区间下发在 M3，跨店汇总在 M8 \\
REQ-008 & 订单状态机 & M5 & 六态转换表驱动，每次变更记来源 \\
REQ-009 & 取消与退款 & M5、M6 & 取消规则在 M5，退款执行在 M6 \\
REQ-010 & 重复提交防护 & M5 & 请求号唯一索引，重放返回原订单号 \\
REQ-011 & 并发抢房 & M4 & 条件更新写进 WHERE 子句，影响 0 行即失败 \\
REQ-012 & 改价与下单隔离 & M3 & 「房型 $\times$ 日期」粒度的锁，只覆盖写价这一个小段 \\
REQ-013 & 库存与可售上限 & M4 & 每晚一行，多晚一次扣减、整单回滚 \\
REQ-014 & 操作日志 & M9 & 只增不改，按角色数据范围过滤查询 \\
REQ-015 & 性能 & 横切 & 见 3.5 决策七 \\
REQ-016 & 安全 & 横切 & 见 3.5 决策八 \\
REQ-017 & 可靠性与数据一致 & 横切 & 见 3.5 决策三、部署图 \\
REQ-018 & 易用性 & 横切 & 入住 3 步，见第 6 章 09 屏 \\
REQ-019 & 范围约束 & 约束 & 不进入本期设计 \\
REQ-020 & 支付超时释放 & M5、M6 & 状态条件更新 + 同事务释放 \\
REQ-021 & 取消截止与前台代取消 & M5 & 18:00 为界，之后转审批 \\
REQ-022 & 未到店标记与迟入住 & M5、M7 & 定时任务判状态，M7 只回答「今天入住了没有」 \\
REQ-023 & 价格快照与定价依据 & M3、M5 & 快照写在订单明细行上，改价不回写 \\
REQ-024 & 房号分配 & M7 & 排房占房记录，唯一约束防重复占用 \\
REQ-025 & 房态变更工单 & M7 & 工单是改房态的唯一入口 \\
REQ-026 & 员工信息可见范围 & M1 & 数据范围规则，店长限本店 \\
REQ-027 & 输入校验与错误提示 & 横切 & 服务端独立校验，见 3.4 节 \\
REQ-028 & 支付与订单一致性 & M6、M5 & 对账补单，库存不足则退款 \\
REQ-029 & 通知失败处理 & M10 & 独立失败队列，不参与订单事务 \\
REQ-030 & 数据保留与删除审批 & M9 & 3 年保留 + 审批后删除 \\
REQ-031 & 订单变更 & M3、M4、M5 & 差额调库存，按新价重算快照 \\
REQ-032 & 多房型多晚次组合 & M4、M5 & 一单多明细，明细粒度是「房型 $\times$ 日期」 \\
REQ-033 & 线上可售上限与线下保留 & M2、M4 & 总数与维修数来自 M2，上限公式在 M4 \\
REQ-034 & 维修中与可售上限联动 & M2、M4、M7 & 维修标记在 M2，逐日扣减在 M4，入口在 M7 \\
REQ-035 & 订单排房时限 & —— & 实验一待确认，本设计不覆盖 \\
REQ-036 & 账号注册与登录 & —— & 实验一待确认，本设计不覆盖 \\
\end{xltabular}
}

\subsection{一笔订单的完整链路}

把上表横过来看：一笔订单从产生到终结，会依次穿过 M2、M3、M4、M5、M6、M7、M8、M10 八个模块。这条链路上有四个「状态会分叉」的时刻，它们是后面所有图与表的重点：

\begin{enumerate}[leftmargin=3.2em]
  \item \textbf{下单}：先扣库存再建单，同一事务。扣库存是链路上唯一会失败的关键动作，放在最前，失败的请求尽早退出，不浪费订单号分配。
  \item \textbf{支付}：回调改状态与超时任务改状态改的是同一行，都带「状态 = 待支付」条件，谁先改成谁生效。
  \item \textbf{入住}：房间从空闲变已预定，订单从已确认变已入住，两个状态机在同一个操作里各走一步。
  \item \textbf{取消}：18:00 前立即生效，18:00 后转审批；未到店订单取消不退房费；重复提交只退一次。
\end{enumerate}
